% Texto: borrador de Laura Segura, con acentos restaurados y citas enlazadas.

\section{Discusión}

Los resultados muestran que es posible operacionalizar principios regulatorios
mediante requerimientos discretos, pruebas computacionales y artefactos
auditables. La principal contribución del marco es que no se limita a declarar
alineación con la Estrategia Nacional de IA, el AI Act o el NIST AI RMF, sino
que produce evidencias verificables: hashes de configuración, registros
estructurados, reportes, model card, datasheet y matriz de cobertura.

La explicabilidad basada en SHAP resultó adecuada para un modelo de árboles y
datos tabulares. Su ventaja principal fue conectar cada inferencia con
atribuciones cuantificadas y, posteriormente, agregar esas atribuciones para
describir el comportamiento global. No obstante, esta selección también impone
una limitación: el desempeño computacional y la validez de las explicaciones
dependen de la familia de modelos y del modo de perturbación seleccionado.

La evaluación de equidad evidenció una restricción importante del dominio. Los
conjuntos públicos de sensores domiciliarios suelen incluir pocos metadatos
demográficos, por lo que no siempre permiten comparar categorías protegidas
como género, etnia, capacidad económica o nivel formativo. El piloto
multi-hogar mitigó parcialmente esta restricción al comparar desempeño entre
viviendas, pero no la elimina por completo. En consecuencia, el marco debe
entenderse como una herramienta para producir evidencia disponible y declarar
sus límites, no como garantía absoluta de ausencia de discriminación.
% Límites que los expedientes declaran y el borrador todavía no menciona:
%   - la franja horaria sale de la hora, que es también variable de entrada:
%     parte de la disparidad entre franjas es por diseño (config.yaml);
%   - la exactitud global del piloto multi-hogar es 0,22--0,43 por vivienda,
%     contra 0,70 en Aruba; un modelo más débil tiene más lugar para mostrar
%     disparidad (D52);
%   - no se midió robustez ni la varianza entre semillas (D56).
\pendiente{decidir si se incluyen los límites que declara el expediente: la hora como variable de entrada, la exactitud baja del multi-hogar, robustez y varianza entre semillas}

La trazabilidad fue el componente transversal que permitió unir resultados
técnicos y responsabilidad. Sin una bitácora persistente, las explicaciones y
los reportes quedarían desconectados de decisiones concretas. Con el esquema
propuesto, cada decisión puede reconstruirse a partir de su identificador,
entrada hasheada, salida, confianza, versión del modelo, explicación y
responsable declarado.

Una diferencia relevante frente a procesos generales de aseguramiento de calidad
en aprendizaje automático, como CRISP-ML(Q), es que el marco no busca optimizar
el modelo ni seleccionar la mejor arquitectura \cite{studer2021crispmlq}. Su
propósito es auditar un sistema dado. Por ello, un veredicto de no cumplimiento
es un resultado válido: revela que el sistema evaluado no satisface los
criterios declarados, pero también demuestra que el procedimiento puede
producir evidencia para sostener esa conclusión.
